\chapter{有理标准形}

在\nameref{sec:实数域上的若当标准形}一节中，我们说明了在实数域这样的非代数闭域上，我们不一定能得到若当标准形. 事实上，我们希望能得到一个在任何数域上都存在的标准形，由于有理数域是最小的数域，因此这种标准形被称为有理标准形. 我们将遵循上一讲\autoref{thm:初等因子组是相似全系不变量} 后提出的三个问题的思路来解决构造有理标准形的问题. 为了阅读方便，我们将这三个问题再次叙述：
\begin{enumerate}
    \item 为什么叫初等因子分解？这是因为这样的分解已经是最细的分解了吗（即不能再进一步分解使其与另一种循环子空间分解同构吗）？
    \item 是否还有其它更粗粒度的分解与其它可能的循环子空间分解对应？即我们是否能结合一些初等因子，得到更大的循环子空间？是否有最粗粒度的分解？
    \item 我们知道在复数域上多项式是可以完全分裂成一次多项式的乘积的，但是在其它一些域，例如实数域或者有理数域上，多项式无法完全分裂，那么将会得到二次多项式或者其它形式的初等因子，此时初等因子分解是否还有意义？
\end{enumerate}

\section{推广的广义特征子空间与第二有理标准形}
为了逻辑上的顺畅，我们首先介绍第二有理标准形的构造，因为这一构造恰好需要对若当标准形的推导中广义特征空间以及初等因子的概念进行推广. 但需要注意的是，有理标准形在复数域上并不能退化为若当标准形，即有理标准形并不是在任何数域上都能找到的一个更广泛的标准形概念，不是若当标准形的推广，这一点是初学时必须要注意的，我们也将在后面详细讲述其中的差异与关联.

回顾\nameref{thm:多项式的唯一分解定理}叙述的唯一分解性质，任意数域上的特征多项式（所以是首一的）$f(\lambda)$都可以被唯一分解为一些首一不可约多项式的乘积，即
\[f(\lambda)=p_1(\lambda)^{n_1}p_2(\lambda)^{n_2}\cdots p_k(\lambda)^{n_k},\]
其中$p_i(\lambda)$是首一不可约多项式，$n_i$是正整数. 如果我们讨论的是复数域，则这些不可约多项式都是一次多项式，即$p_i(\lambda)=\lambda-\lambda_i$，其中$\lambda_i$是复数域上的特征值. 但是如果我们讨论的是实数域或者有理数域，则这些不可约多项式可能是一次多项式，也可能是更高次的多项式. 事实上，在复数域上我们推导若当标准形的起点是定义广义特征子空间，然后将广义特征子空间分解为循环子空间的直和，那么我们最直接的想法便是，在所有数域上（包括非代数闭域）定义类似的概念，进行类似的分解. 于是我们首先有如下定义：
\begin{definition}{}{推广广义特征子空间}
    设$V$是一个有限维线性空间，$T\in\mathcal{L}(V)$. 设$T$的特征多项式为
    \[f(\lambda)=p_1(\lambda)^{n_1}p_2(\lambda)^{n_2}\cdots p_k(\lambda)^{n_k},\]
    其中$p_i(\lambda)$是首一不可约多项式，$n_i$是正整数. 对于每一个不可约多项式$p_i(\lambda)$，我们定义\term{广义特征子空间}$G_{p_i}$为
    \[G_{p_i}=\{v\in V\mid\exists m\in\mathbf{N}^+,(p_i(T))^mv=0\},\]
    $G_{p_i}$中的向量则称为\term{广义特征向量}.
\end{definition}

事实上我们很容易发现这就是原先复数域上定义的广义特征子空间在任意数域上的推广形式：当数域为复数域时，$p_i(\lambda)=\lambda-\lambda_i$，其中$\lambda_i$是复数域上的特征值，则$G_{p_i}=\{v\in V\mid\exists m\in\mathbf{N}^+,(T-\lambda_iI)^mv=0\}$，实际上就等同于之前定义的广义特征子空间. 于是，我们也期望将\autoref{thm:广义特征性质} 推广到任意数域上：
\begin{theorem}{}{推广广义特征性质}
    设$V$是有限维线性空间，$T\in \mathcal{L}(V)$. $T$的特征多项式有如下分解：
    \[f(\lambda)=p_1(\lambda)^{n_1}p_2(\lambda)^{n_2}\cdots p_k(\lambda)^{n_k},\]
    其中$p_i(\lambda)$是首一不可约多项式，$n_i$是正整数，$G_{p_i}$为广义特征子空间，则我们有如下结论：
    \begin{enumerate}[label=(\arabic*)]
        \item 每个$G_{p_i}$在$T$下都是不变的；
        \item $T$对应于不同不可约多项式的广义特征向量线性无关；
        \item $T$对应于不同不可约多项式的广义特征子空间的和为直和，且$V=G_{p_1}\oplus G_{p_2}\oplus\cdots\oplus G_{p_k}$；
        \item $V$有一个由$T$的广义特征向量组成的基.
    \end{enumerate}
\end{theorem}
\begin{proof}
    \begin{enumerate}
        \item 设$v\in G_{p_i}$，则存在正整数$m$使得$(p_i(T))^mv=0$. 由于$p_i(T)$与$T$可交换，因此
              \[
                  (p_i(T))^m(Tv)=T(p_i(T))^mv=0.
              \]
              所以$Tv\in G_{p_i}$，即$G_{p_i}$在$T$下不变.

        \item 我们先证明第(3)点，再由第(3)点推出本结论.

        \item 记
              \[
                  f_i(\lambda)=\frac{f(\lambda)}{p_i(\lambda)^{n_i}}=\prod_{j\neq i}p_j(\lambda)^{n_j}.
              \]
              由于$p_i^{n_i}$与$f_i$互素，若$v\in G_{p_i}$，设$(p_i(T))^mv=0$，则$p_i^m$与$f_i$也互素. 根据裴蜀定理，存在$a,b\in\mathbf{F}[x]$使得
              \[
                  a(\lambda)p_i(\lambda)^m+b(\lambda)f_i(\lambda)=1.
              \]
              代入$T$后作用于$v$，得到
              \[
                  v=a(T)(p_i(T))^mv+b(T)f_i(T)v=b(T)f_i(T)v.
              \]
              再作用$(p_i(T))^{n_i}$，便有
              \[
                  (p_i(T))^{n_i}v=b(T)f(T)v=0,
              \]
              因为由 Hamilton-Cayley 定理知$f(T)=0$. 这说明$G_{p_i}\subseteq \ker (p_i(T))^{n_i}$；反过来，若$v\in \ker (p_i(T))^{n_i}$，则它当然属于$G_{p_i}$. 因而
              \[
                  G_{p_i}=\ker (p_i(T))^{n_i}.
              \]

              现在由$f(T)=0$可知$\ker f(T)=V$，再利用\autoref{cor:多项式分解与核空间直和2}，得到
              \[
                  V=\ker f(T)=\ker (p_1(T))^{n_1}\oplus\cdots\oplus\ker (p_k(T))^{n_k}=G_{p_1}\oplus\cdots\oplus G_{p_k}.
              \]
              因此第(3)点成立.

        \item 由第(3)点，任取每个$G_{p_i}$的一组基，将它们合并便得到$V$的一组基. 而$G_{p_i}$中的每个向量按定义都是$T$对应于$p_i$的广义特征向量，因此这组基由广义特征向量组成.
    \end{enumerate}

    第(2)点现在是第(3)点的直接推论：若$v_1,\ldots,v_s$分别属于互不相同的广义特征子空间$G_{p_{i_1}},\ldots,G_{p_{i_s}}$且都非零，而有
    \[
        c_1v_1+\cdots+c_sv_s=0,
    \]
    那么直和分解的唯一性迫使$c_1v_1=\cdots=c_sv_s=0$，于是$c_1=\cdots=c_s=0$. 故它们线性无关.
\end{proof}

接下来我们的目标便是将广义特征子空间分解为循环子空间的直和. 事实上，我们无法分解为若当块所需的循环子空间的形式，因为一方面我们不一定在复数域上讨论，另一方面若当标准形有唯一性，而我们需要找到另一种标准形，如果还保持原先的循环子空间还是只能局限于若当标准形. 因此我们使用\autoref{def:循环子空间}中定义的循环子空间来进行分解.

若我们记由$v$生成的$T-\text{循环子空间}$为$C_v$，则我们可以得到限制映射$T|_{C_v}$在循环基$\{v,Tv,\ldots,T^{m-1}v\}$下的矩阵. 为了得到这一矩阵，我们设$Tv=a_0v+a_1Tv+\cdots+a_{m-1}T^{m-1}v$（因为$Tv$一定可以在这组基下被表示），则我们有矩阵表示为
\[A=\begin{pmatrix}
        0      & 0      & \cdots & 0      & -a_0     \\
        1      & 0      & \cdots & 0      & -a_1     \\
        0      & 1      & \cdots & 0      & -a_2     \\
        \vdots & \vdots & \ddots & \vdots & \vdots   \\
        0      & 0      & \cdots & 1      & -a_{m-1}
    \end{pmatrix},\]
事实上我们可以很容易计算出这一矩阵的特征多项式为$|\lambda E-A|=\lambda^m+a_{m-1}\lambda^{m-1}+\cdots+a_1\lambda+a_0$. 更进一步地，根据\autoref{thm:特征多项式等于极小多项式}，这一矩阵的极小多项式就等于这一特征多项式. 一件很有趣的事情是，这一矩阵除了主对角线下一行的元素全为1之外，只有最后一列元素非零，且是特征多项式的负系数的排列，因此这一矩阵被称为多项式$\lambda^m+a_{m-1}\lambda^{m-1}+\cdots+a_1\lambda+a_0$的\term{友阵}，也称一个\term{弗罗贝尼乌斯块}. 如果一个分块对角矩阵的对角块都是形如上述的弗罗贝尼乌斯块，则这一矩阵被称为\term{弗罗贝尼乌斯标准形}，也即\term{有理标准形}.

于是接下来的目标则是将广义特征子空间分解为上述可以生成弗罗贝尼乌斯块的循环子空间的直和. 接下来的几个定理将逐步实现这一目标：
\begin{theorem}{}{}
    若$T$的极小多项式具有如下形式：$g(\lambda)=(p(\lambda))^m$，则$T$有一组基使得其表示矩阵为有理标准形.
\end{theorem}
\begin{proof}
    记$d=\deg p$. 我们实际上证明一个更强的结论：$V$可以分解为若干个$T$-循环子空间的直和，
    \[
        V=C_{v_1}\oplus\cdots\oplus C_{v_s},
    \]
    且每个$C_{v_i}$上的极小多项式都形如$p^{r_i}$，并且向量$p(T)^{r_i-1}v_i$在$\ker p(T)$中构成$E=\mathbf{F}[x]/(p)$上的一组基. 取每个$C_{v_i}$的循环基并起来，便得到所需的有理标准形基.

    我们对$m$作归纳.

    当$m=1$时，$p(T)=0$. 记$E=\mathbf{F}[x]/(p)$. 由于$p$不可约，$E$是一个域. 对任意$\overline{f(x)}\in E$与$v\in V$，定义
    \[
        \overline{f(x)}\cdot v=f(T)v.
    \]
    这是良定义的，因为若$f-g$能被$p$整除，则$(f-g)(T)=0$. 因而$V$成为$E$上的向量空间. 取$V$的一组$E$-基$u_1,\ldots,u_s$. 对每个$u_i\neq 0$，由于$p(T)u_i=0$，其极小多项式必为$p$，于是$C_{u_i}$是一个$d$维循环子空间，其循环基可取为
    \[
        u_i,Tu_i,\ldots,T^{d-1}u_i.
    \]
    将这些循环基合并便得到$V$的一组$\mathbf{F}$-基，而且$u_1,\ldots,u_s$本身就是$\ker p(T)=V$的一组$E$-基，故结论成立.

    设$m>1$，并假设命题对$p^{m-1}$成立. 令
    \[
        U=\im p(T).
    \]
    由于$p(T)$与$T$可交换，$U$在$T$下不变. 又$T\vert_U$的极小多项式整除$p^{m-1}$，故由归纳假设，$U$可以分解为若干个循环子空间的直和
    \[
        U=C_{w_1}\oplus\cdots\oplus C_{w_s},
    \]
    其中每个$C_{w_i}$上的极小多项式为$p^{r_i-1}$，并且$p(T)^{r_i-2}w_i$在$\ker p(T)\cap U$中构成$E=\mathbf{F}[x]/(p)$上的一组基.

    对每个$i$，取$v_i\in V$使得$p(T)v_i=w_i$. 我们断言$C_{v_i}$上的极小多项式为$p^{r_i}$. 事实上，
    \[
        p(T)^{r_i}v_i=p(T)^{r_i-1}w_i=0.
    \]
    若$p(T)^{r_i-1}v_i=0$，则
    \[
        p(T)^{r_i-2}w_i=p(T)^{r_i-1}v_i=0,
    \]
    这与$C_{w_i}$上的极小多项式是$p^{r_i-1}$矛盾. 因此$C_{v_i}$上的极小多项式确为$p^{r_i}$.

    由
    \[
        p(T)^{r_i-1}v_i=p(T)^{r_i-2}w_i
    \]
    可知这些向量在$\ker p(T)$中线性无关，并在$\ker p(T)\cap U$中构成一组$E$-基. 将它们扩充为$\ker p(T)$的一组$E$-基，不妨记扩充得到的向量为$u_1,\ldots,u_t$. 由于每个$u_j\neq 0$且$p(T)u_j=0$，所以$C_{u_j}$上的极小多项式都是$p$.

    现在定义
    \[
        W=C_{v_1}+\cdots+C_{v_s}+C_{u_1}+\cdots+C_{u_t}.
    \]
    显然$W$在$T$下不变. 并且
    \[
        p(T)W=C_{w_1}+\cdots+C_{w_s}=U,
    \]
    因为$p(T)C_{v_i}=C_{w_i}$，而$p(T)C_{u_j}=0$. 此外，根据$u_1,\ldots,u_t$的选取方式，$\ker p(T)\subseteq W$. 故对限制映射$p(T)\vert_W$应用秩-零度定理，有
    \[
        \dim W=r\bigl(p(T)\vert_W\bigr)+\dim\ker\bigl(p(T)\vert_W\bigr)=\dim U+\dim\ker p(T).
    \]
    再对$p(T):V\to V$应用秩-零度定理，得到
    \[
        \dim V=\dim U+\dim\ker p(T).
    \]
    从而$\dim W=\dim V$，所以$W=V$.

    最后证明上述和是直和. 对每个$i$，$C_{v_i}$的极小多项式是$p^{r_i}$，故$\dim C_{v_i}=dr_i$；而$\dim C_{w_i}=d(r_i-1)$，所以
    \[
        \dim C_{v_i}=\dim C_{w_i}+d.
    \]
    又每个$C_{u_j}$的维数都是$d$. 因此
    \[
        \sum_{i=1}^s\dim C_{v_i}+\sum_{j=1}^t\dim C_{u_j}
        =\sum_{i=1}^s\bigl(\dim C_{w_i}+d\bigr)+td
        =\dim U+\dim\ker p(T)
        =\dim V.
    \]
    而这些子空间的和已经等于$V$，故该和必为直和.

    于是$V$被分解为若干个循环子空间的直和. 对每个循环子空间取循环基，$T$在这组基下的矩阵就是若干个友阵构成的分块对角矩阵，即有理标准形.
\end{proof}

\begin{corollary}{}{}
    设$p$是$T$的特征多项式的一个首一不可约因式，则$G_p$有一组由不相交的循环基的并集构成的基.
\end{corollary}
\begin{proof}
    由\autoref{thm:推广广义特征性质}，$G_p$在$T$下不变. 设$T\vert_{G_p}$的极小多项式为$p^m$，则它确实是某个$p$的幂，因为其极小多项式整除$T\vert_{G_p}$的特征多项式，而后者是$p$的幂. 因此可将上一定理应用于$T\vert_{G_p}$，便得到$G_p$有一组由不相交的循环基的并集构成的基.
\end{proof}

\begin{corollary}{}{第二有理基存在}
    设$V$是$\mathbf{F}$上的有限维线性空间，$T\in\mathcal{L}(V)$，则存在一组基$B$使得$T$在$B$下的矩阵表示为有理标准形.
\end{corollary}
\begin{proof}
    根据\autoref{thm:推广广义特征性质}，若$T$的特征多项式分解为
    \[
        f(\lambda)=p_1(\lambda)^{n_1}\cdots p_k(\lambda)^{n_k},
    \]
    则
    \[
        V=G_{p_1}\oplus\cdots\oplus G_{p_k},
    \]
    且每个$G_{p_i}$都在$T$下不变. 再由上一个推论可知，每个$G_{p_i}$都有一组由不相交的循环基的并集构成的基. 将这些基合并，就得到$V$的一组基，使得$T$在这组基下的矩阵表示为若干个友阵构成的分块对角矩阵，也就是有理标准形.
\end{proof}
事实上，如果我们按照上一讲的做法将直和分解书写出来，我们发现这样的分解仍然是基于初等因子的分解，我们将这样得到的有理标准形称为\term{第二有理标准形}.

接下来我们需要描述矩阵的第二有理标准形.
\begin{corollary}{}{}
    设$A\in\mathbf{M}_n(\mathbf{F})$，则存在可逆矩阵$P$，使得$P^{-1}AP$是分块对角矩阵，且每个对角块都是某个首一不可约多项式的幂的友阵. 这一矩阵称为$A$的\term{第二有理标准形}.
\end{corollary}
\begin{proof}
    将$A$看作线性变换$\sigma:\mathbf{F}^n\to\mathbf{F}^n,\ \sigma(x)=Ax$在自然基下的矩阵表示. 由\autoref{cor:第二有理基存在}，存在$\mathbf{F}^n$的一组基$B$使得$\sigma$在$B$下的矩阵表示为有理标准形. 设$P$为从自然基变到$B$的过渡矩阵，则由换基公式可知$P^{-1}AP$就是这一个有理标准形.
\end{proof}

有理标准形的唯一性定理：
\begin{theorem}{}{第二有理标准形唯一}
    设$p$是$T$的特征多项式的一个首一不可约因式，$d=\deg p$. 在$T$的第二有理标准形中，设所有以$p$为底的初等因子为
    \[
        p^{m_1},p^{m_2},\ldots,p^{m_s}\qquad (m_1\geqslant m_2\geqslant\cdots\geqslant m_s).
    \]
    对任意正整数$i$，记$r_i$为满足$m_j\geqslant i$的指标$j$的个数，即$r_i$表示指数至少为$i$的$p$-初等因子的个数，则
    \begin{enumerate}
        \item $r_1=\dfrac{1}{d}(\dim V-r(p(T)))$；
        \item $r_i=\dfrac{1}{d}(r(p(T)^{i-1})-r(p(T)^i))\quad (i>1)$.
    \end{enumerate}
    从而第二有理标准形除对角块的排列次序外唯一.
\end{theorem}
\begin{proof}
    设$V$已经写成第二有理标准形对应的循环子空间直和. 对于每一个以$p$为底的初等因子$p^{m_j}$，对应的循环子空间同构于$\mathbf{F}[x]/(p^{m_j})$. 在这个商空间上，$p(T)$对应于“乘以$p(x)$”这一算子，因此对任意$i\geqslant 1$，
    \[
        \dim\ker\bigl(p(T)^i\vert_{\mathbf{F}[x]/(p^{m_j})}\bigr)=d\min\{i,m_j\}.
    \]
    于是当$i=1$时，每个这样的循环子空间都向$\ker p(T)$贡献恰好$d$维；当$i>1$时，它向$\ker p(T)^i$相对于$\ker p(T)^{i-1}$的增长贡献恰好是
    \[
        \dim\ker p(T)^i-\dim\ker p(T)^{i-1}=
        \begin{cases}
            d, & m_j\geqslant i, \\
            0, & m_j<i.
        \end{cases}
    \]
    另一方面，若某个循环子空间对应的初等因子底数不是$p$，则由于该底数与$p$互素，$p(T)$在这个循环子空间上可逆，因此对上述各个核空间都没有贡献.

    故有
    \[
        \dim\ker p(T)=dr_1,
    \]
    以及对$i>1$，
    \[
        \dim\ker p(T)^i-\dim\ker p(T)^{i-1}=dr_i.
    \]
    再由秩-零度定理，
    \[
        \dim\ker p(T)=\dim V-r(p(T)),
    \]
    且
    \[
        \dim\ker p(T)^i-\dim\ker p(T)^{i-1}=r(p(T)^{i-1})-r(p(T)^i).
    \]
    代入即得题中的两个公式.

    最后，大小恰为$i$的$p$-初等因子的个数等于$r_i-r_{i+1}$，故所有$m_j$都被唯一确定. 对每个首一不可约多项式$p$都作同样讨论，便可唯一确定全部初等因子，因此第二有理标准形除对角块排列次序外唯一.
\end{proof}

\section{不变因子与第一有理标准形}
接下来我们将讨论第二个问题及其回答，从而引入不变因子的概念以及第一有理标准形. 事实上，尽管我们暂未给出第一个问题的解答，但直觉告诉我们，初等因子分解已经是一种比较细的分解方式了——因为它都是一些准素多项式，因此我们希望对一些初等因子进行合并，得到更大的循环子空间. 因此我们需要研究初等因子合并的条件，我们将其叙述为以下定理：
\begin{theorem}{}{初等因子组合}
    设$p_1(\lambda),p_2(\lambda)\in\mathbf{F}[\lambda]$，则循环子空间$C_1\cong\mathbf{F}[\lambda]/(p_1(\lambda))$与$C_2\cong\mathbf{F}[\lambda]/(p_2(\lambda))$可以直和成一个新的循环子空间$C\cong\mathbf{F}[\lambda]/(p(\lambda))$（其中$p(\lambda)\in\mathbf{F}[\lambda]$）当且仅当$p_1(\lambda)$与$p_2(\lambda)$互素，且互素时有$p(\lambda)=p_1(\lambda)p_2(\lambda)$.
\end{theorem}
\begin{proof}
    先证充分性. 若$p_1(\lambda)$与$p_2(\lambda)$互素，设$p(\lambda)=p_1(\lambda)p_2(\lambda)$，并定义线性映射
    \[
        \Phi:\mathbf{F}[\lambda]/(p(\lambda))\to \mathbf{F}[\lambda]/(p_1(\lambda))\oplus\mathbf{F}[\lambda]/(p_2(\lambda)),
        \qquad
        \overline{f(\lambda)}\mapsto \bigl(\overline{f(\lambda)},\overline{f(\lambda)}\bigr).
    \]
    若$\Phi(\overline{f(\lambda)})=0$，则$p_1(\lambda)\mid f(\lambda)$且$p_2(\lambda)\mid f(\lambda)$. 由于$p_1,p_2$互素，故$p_1(\lambda)p_2(\lambda)\mid f(\lambda)$，即$\overline{f(\lambda)}=0$，所以$\Phi$单射. 又
    \[
        \dim \mathbf{F}[\lambda]/(p(\lambda))=\deg p=\deg p_1+\deg p_2
    \]
    与
    \[
        \dim \bigl(\mathbf{F}[\lambda]/(p_1(\lambda))\oplus\mathbf{F}[\lambda]/(p_2(\lambda))\bigr)=\deg p_1+\deg p_2
    \]
    相同，因此$\Phi$是同构. 这就说明
    \[
        \mathbf{F}[\lambda]/(p_1(\lambda)p_2(\lambda))\cong \mathbf{F}[\lambda]/(p_1(\lambda))\oplus\mathbf{F}[\lambda]/(p_2(\lambda)),
    \]
    即$C_1$与$C_2$可以直和成一个新的循环子空间.

    再证必要性. 设存在首一多项式$p(\lambda)$使得
    \[
        \mathbf{F}[\lambda]/(p(\lambda))\cong \mathbf{F}[\lambda]/(p_1(\lambda))\oplus\mathbf{F}[\lambda]/(p_2(\lambda)).
    \]
    由于右侧空间的维数为$\deg p_1+\deg p_2$，故$\deg p=\deg p_1+\deg p_2$.

    另一方面，$p_1(\lambda)$与$p_2(\lambda)$分别消去右侧的两个直和项，所以它们的最小公倍式$l(\lambda)=\lcm(p_1(\lambda),p_2(\lambda))$消去整个直和空间. 因而$l(\lambda)$也是$\mathbf{F}[\lambda]/(p(\lambda))$上的零化多项式. 由循环子空间的极小多项式等于特征多项式，知道$p(\lambda)$正是该空间的极小多项式，所以$p(\lambda)\mid l(\lambda)$. 从而
    \[
        \deg p\leqslant \deg l\leqslant \deg p_1+\deg p_2=\deg p.
    \]
    因此$\deg l=\deg p_1+\deg p_2$. 这只可能在$p_1(\lambda)$与$p_2(\lambda)$互素时发生；否则$l(\lambda)$的次数将严格小于$\deg p_1+\deg p_2$. 故$p_1(\lambda)$与$p_2(\lambda)$必互素. 此时
    \[
        l(\lambda)=p_1(\lambda)p_2(\lambda),
    \]
    又$p(\lambda)$是首一且与$l(\lambda)$次数相同，并满足$p(\lambda)\mid l(\lambda)$，故$p(\lambda)=p_1(\lambda)p_2(\lambda)$.
\end{proof}

\begin{corollary}{}{}
    设$p_1(\lambda),\ldots,p_s(\lambda)\in\mathbf{F}[\lambda]$都是首一多项式，则循环子空间
    \[
        C_i\cong \mathbf{F}[\lambda]/(p_i(\lambda))\qquad (i=1,\ldots,s)
    \]
    可以直和成一个新的循环子空间，当且仅当$p_1(\lambda),\ldots,p_s(\lambda)$两两互素；此时所得循环子空间同构于
    \[
        \mathbf{F}[\lambda]/(p_1(\lambda)\cdots p_s(\lambda)).
    \]
\end{corollary}
\begin{proof}
    对$s$作归纳即可. 当$s=2$时结论就是上一个定理. 设结论对$s-1$个多项式成立，则前$s-1$个对应的循环子空间可组合成一个循环子空间，当且仅当它们两两互素；再将这个新的循环子空间与第$s$个循环子空间应用上一个定理，便得到结论.
\end{proof}

事实上一个巧合在于这一定理间接回答了第一个问题：
\begin{corollary}{}{}
    \autoref{eq:19:初等因子分解} 给出的初等因子分解是循环子空间分解中最细的一种：其中每一个直和项都不能再分解为两个非平凡循环子空间的直和，并且不同初等因子之间也不能再作进一步的组合.
\end{corollary}
\begin{proof}
    对于\autoref{eq:19:初等因子分解} 中的任意一个直和项$\mathbf{F}[\lambda]/(p(\lambda)^r)$，若它还能分解为两个非平凡循环子空间的直和，则根据\autoref{thm:初等因子组合}，$p(\lambda)^r$应能写成两个互素首一多项式的乘积. 但$p(\lambda)^r$的任意两个非平凡因式都含有公因子$p(\lambda)$，不可能互素，故这是不可能的.

    另一方面，若想把两个不同的初等因子继续组合成更大的循环子空间，根据\autoref{thm:初等因子组合}，它们对应的多项式必须互素. 但在初等因子分解中，凡是可以组合的不同底数的初等因子都已经由上一个推论统一控制，而同一底数的两个初等因子显然不互素，不能组合. 故初等因子分解已经是最细分解.
\end{proof}

自然地，在证明初等因子分解是最细分解后，我们可以利用\autoref{thm:初等因子组合} 来构造更大的循环子空间. 事实上，根据定理，我们知道这样的构造由非常多种，我们可以仅仅合并两个初等因子，也可以把所有可以合并的初等因子全部合并. 为了得到更标准且优雅的形式，我们做如下合并：首先我们将$T$的特征多项式进行分解得到
\[f(\lambda)=p_1(\lambda)^{n_1}p_2(\lambda)^{n_2}\cdots p_k(\lambda)^{n_k},\]
其中$p_i(\lambda)$是首一不可约多项式，$n_i$是正整数. 于是$T$的初等因子具有形式$p_i^{r_{ij}}$，我们将这些初等因子按照不可约多项式进行分组，即将所有底数$p_i(\lambda)$相同的初等因子写在一行，并且按幂次降序排列：
\begin{gather*}
    p_1^{r_{11}},p_1^{r_{12}},\ldots,p_1^{r_{1s_1}}, \\
    p_2^{r_{21}},p_2^{r_{22}},\ldots,p_2^{r_{2s_2}}, \\
    \cdots \\
    p_k^{r_{k1}},p_k^{r_{k2}},\ldots,p_k^{r_{ks_k}}.
\end{gather*}
其中$r_{ij}$是正整数，$r_{i1}\geqslant r_{i2}\geqslant\cdots\geqslant r_{is_i}$，且有$\sum\limits_{j=1}^{s_i}r_{ij}=n_i$. 接下来我们将每一行的长度对齐，不足的补$1$，于是可以写成
\begin{gather*}
    p_1^{r_{11}},p_1^{r_{12}},\ldots,p_1^{r_{1s}}, \\
    p_2^{r_{21}},p_2^{r_{22}},\ldots,p_2^{r_{2s}}, \\
    \cdots \\
    p_k^{r_{k1}},p_k^{r_{k2}},\ldots,p_k^{r_{ks}}.
\end{gather*}
其中$s=\max\{s_1,s_2,\ldots,s_k\}$. 然后我们从后往前依次计算每一列的元素相乘的结果，即
\begin{align*}
    d_1     & =p_1^{r_{1s}}p_2^{r_{2s}}\cdots p_k^{r_{ks}}, \\
            & \cdots                                        \\
    d_{s-1} & =p_1^{r_{12}}p_2^{r_{22}}\cdots p_k^{r_{k2}}, \\
    d_s     & =p_1^{r_{11}}p_2^{r_{21}}\cdots p_k^{r_{k1}}.
\end{align*}
根据\autoref{thm:初等因子组合}，由于$p_i(\lambda)$两两互素，因此我们有
\[\mathbf{F}[x]/(d_{s+1-i})=\mathbf{F}[x]/(p_1^{r_{1i}})\oplus\mathbf{F}[x]/(p_2^{r_{2i}})\oplus\cdots\oplus\mathbf{F}[x]/(p_k^{r_{ki}}),\]
且$\mathbf{F}[x]/(d_i)$可以由一组循环基张成. 因此\autoref{eq:19:初等因子分解} 可以写成经过上述合并的形式：
\begin{equation} \label{eq:20:不变因子分解}
    V=\mathbf{F}[x]/(d_1)\oplus\mathbf{F}[x]/(d_2)\oplus\cdots\oplus\mathbf{F}[x]/(d_s),
\end{equation}
其中$d_i|d_{i+1}(i=1,\ldots,s-1)$. 这样的一组具有``一个整除下一个''性质的多项式被称为$T$的\term{不变因子组}，其中的每个多项式$d_i$都是$T$的一个\term{不变因子}.

事实上，由\autoref{thm:初等因子组合} 的互素条件可知，这一分解已经是最``粗粒度''的分解了，因为每个不变因子都是不互素的，无法进一步组合成更大的循环子空间. 接下来我们需要为这一不变因子分解具体实现出一组循环基和对应的标准形.

\begin{theorem}{}{第一有理基构造}
    在\autoref{eq:20:不变因子分解} 的记号下，存在一组基$B$使得$T$在$B$下的矩阵表示为
    \[
        \diag(C(d_1),C(d_2),\ldots,C(d_s)),
    \]
    其中$C(d_i)$表示多项式$d_i$的友阵. 这一分块对角矩阵称为$T$的\term{第一有理标准形}.
\end{theorem}
\begin{proof}
    对于\autoref{eq:20:不变因子分解} 中的每一个直和项$\mathbf{F}[x]/(d_i)$，考虑其自然生成元$\overline{1}$所生成的循环子空间. 由\autoref{thm:循环子空间}，它的一组循环基可取为
    \[
        \overline{1},\overline{x},\ldots,\overline{x^{\deg d_i-1}}.
    \]
    在这组基下，“乘以$x$”这一线性映射的矩阵恰好就是$d_i$的友阵$C(d_i)$. 因此每一个直和项都能给出一个友阵块.

    现在将各个直和项上的上述循环基依次合并，便得到整个空间$V$的一组基. 因为\autoref{eq:20:不变因子分解} 是直和分解，所以$T$在这组基下的矩阵表示就是
    \[
        \diag(C(d_1),C(d_2),\ldots,C(d_s)).
    \]
    这正是第一有理标准形.
\end{proof}

\begin{theorem}{}{}
    不变因子组也是相似的全系不变量.
\end{theorem}
\begin{proof}
    由上一节的\autoref{thm:第二有理标准形唯一} 可知，初等因子组是相似的全系不变量. 因此只需证明初等因子组与不变因子组可以互相唯一确定.

    由初等因子组到不变因子组的确定方式，就是\autoref{eq:20:不变因子分解} 前面的“按底数分行、按幂次对齐、按列相乘”的过程，这一过程显然是唯一的.

    反过来，已知不变因子组$d_1,\ldots,d_s$，把每个$d_i$在$\mathbf{F}[x]$中唯一分解为首一不可约多项式的幂的乘积，然后把相同底数的幂次按列重新排开，便能唯一恢复出全部初等因子. 也就是说，不变因子组与初等因子组一一对应.

    于是两个矩阵相似当且仅当它们有相同的不变因子组，即不变因子组也是相似的全系不变量.
\end{proof}

\begin{corollary}{}{}
    第一有理标准形除对角块的排列次序外唯一.
\end{corollary}

\begin{summary}

\end{summary}

\begin{exercise}
    % \exquote[]{}

    \begin{exgroup}
        \item
    \end{exgroup}

    \begin{exgroup}
        \item
    \end{exgroup}

    \begin{exgroup}
        \item
    \end{exgroup}
\end{exercise}
